\label{lecture02}
\subsection*{Цели и маршрут занятия (90 минут)}
После лекции студент должен уметь отличать значение от его кода, оценивать
объём несжатых данных, объяснять представление целых и вещественных чисел
и читать простые разрядные схемы. Главная идея занятия:
\begin{quote}
Биты сами по себе не являются ни текстом, ни числом, ни командой.
Значение определяется форматом и правилами интерпретации.
\end{quote}
Введение занимает 5 минут, графика --- 10, звук --- 6, текст --- 9,
целые числа --- 25, IEEE 754 и точность --- 20, машинные команды,
маски и сдвиги --- 10, итоги --- 5 минут. Последний блок опирается
на побитовые операции, введённые при рассмотрении целых чисел.
Подробная арифметика с плавающей точкой вынесена в дополнительный материал
и углублённо разбирается на занятии №3.

\subsection{Битовая последовательность и её смысл}
Бит принимает одно из двух значений: 0 или 1. В рассматриваемых машинах
байт содержит 8 бит. Из $n$ бит можно составить $2^n$ различных комбинаций.
Чтобы пронумеровать $N$ объектов кодами одинаковой длины, необходимо
$n=\lceil\log_2 N\rceil$ бит (для $N=1$ математически достаточно нуля бит).
Например, для 5 состояний нужно 3 бита: двух бит хватило бы только на 4.

\textbf{Формат} определяет длины и назначение полей, а \textbf{правило
кодирования} связывает значения с комбинациями битов. Байт
\texttt{01000001} может означать число 65 или латинскую букву A в ASCII.
Байт \texttt{11111111} означает 255 как беззнаковое число и $-1$ в
8-битном дополнительном коде. Адрес в памяти сообщает, \emph{где} лежат
байты, но не задаёт сам по себе их смысл.

\textbf{Единицы объёма:} 1 kB = 1000 байт, 1 MB = $10^6$ байт;
1 KiB = $2^{10}=1024$ байта, 1 MiB = $2^{20}$ байт.
В задачах всегда указываем, какие единицы используем.

\subsection{Кодирование графической информации}
\subsubsection*{Растр, палитра и RGB}
\textbf{Растровое изображение} --- прямоугольная сетка пикселей.
Размер $W\times H$ означает $W$ столбцов и $H$ строк.
В простом чёрно-белом изображении без полутонов достаточно 1 бита
на пиксель. Соответствие 0 и 1 цветам задаёт формат: оно не универсально.

При \textbf{палитровом} кодировании в пикселе хранится номер цвета,
а таблица цветов хранится отдельно. Если на номер отведено $p$ бит,
можно выбрать до $2^p$ цветов. Четыре цвета требуют 2 бит на индекс,
16 цветов --- 4 бит. Состав палитры может быть произвольным.

В \textbf{RGB} цвет задаётся интенсивностями красного, зелёного и синего.
Для RGB888 на канал отведено 8 бит: каждый принимает значения от 0 до 255.
На пиксель приходится 24 бита, всего $2^{24}=16\,777\,216$ комбинаций.
Например, $(255,0,0)$ --- красный цвет, $(255,255,0)$ --- жёлтый.
Чёрному соответствует $(0,0,0)$, белому --- $(255,255,255)$.
Если каждому каналу дать только 1 бит, получим 8 комбинаций RGB;
сочетание $(1,1,0)$ также даёт жёлтый, а не коричневый.
Формат RGBA8888 дополнительно хранит 8-битный альфа-канал прозрачности
и занимает 32 бита на пиксель.

\subsubsection*{Оценка объёма}
При плотной упаковке без заголовков, палитры, выравнивания строк и сжатия:
\[
V_{\text{бит}}=W H p,\qquad
V_{\text{байт}}=\left\lceil\frac{W H p}{8}\right\rceil.
\]
\textit{Пример:} $640\times480$ пикселей в RGB888 занимают
$640\cdot480\cdot3=921\,600$ байт, то есть 900 KiB.
В палитровом формате с 16 цветами сами индексы занимают
$640\cdot480\cdot4/8=153\,600$ байт, или 150 KiB;
палитру нужно учитывать дополнительно.
Размер PNG или JPEG по этой формуле определить нельзя: он зависит
от содержимого, сжатия и служебных данных.

\subsubsection*{Векторное представление}
Векторный файл описывает объекты: отрезки, кривые, заливки, текст,
их координаты и свойства. При выводе на экран объекты растеризуются.
Например, SVG может хранить круг через координаты центра и радиус,
а не через цвета всех пикселей. Вектор удобен для схем и логотипов,
растр --- для фотографий. Количество пикселей влияет на детализацию
растра, но не является единственным показателем качества изображения.

\subsection{Кодирование звука}
Для цифровой записи аналоговый сигнал измеряют через равные промежутки
времени. Это \textbf{дискретизация}. Частота $f_s$ показывает число
отсчётов в секунду в каждом канале. При $f_s=48\,000$ Гц интервал
между отсчётами равен $1/48\,000$ секунды.

Измеренное значение округляют до одного из допустимых уровней --- это
\textbf{квантование}. При $b$ битах на отсчёт доступны $2^b$ кодов.
Разрядность влияет на ошибку квантования, а частота дискретизации ---
на представимые частоты сигнала. Это разные характеристики.
Для восстановления ограниченного по частоте сигнала частота дискретизации
должна быть больше удвоенной наивысшей частоты; перед оцифровкой
используют фильтр против наложения спектров.

В несжатом PCM при $c$ каналах и длительности $t$ секунд:
\[
V_{\text{бит}}=f_s b c t.
\]
\textit{Пример:} 10 секунд, 48 kHz, 16 бит, стерео:
$48\,000\cdot16\cdot2\cdot10/8=1\,920\,000$ байт, или 1.92 MB.
Стерео означает два канала, а не удвоенную частоту каждого канала.
Заголовок контейнера и сжатие здесь не учитываются.

\subsection{Кодирование текста}
\subsubsection*{ASCII и однобайтные кодировки}
ASCII --- 7-битный стандарт со 128 кодами от 0 до 127.
Он включает латинские буквы, цифры, знаки и управляющие коды.
Например, A имеет код 65 (\texttt{0x41}), символ 0 --- код 48
(\texttt{0x30}), перевод строки LF --- код 10 (\texttt{0x0A}).
\emph{Число 0 и символ 0 --- разные значения и представления.}

В байте доступны 256 кодов, но единой «второй половины ASCII» нет.
Windows-1251, CP866 и другие однобайтные кодировки по-разному назначают
коды 128--255. Поэтому одни и те же байты, прочитанные в неверной кодировке,
могут дать нечитаемый текст.

\subsubsection*{Unicode и формы кодирования}
Unicode назначает символам \textbf{кодовые точки}, например A ---
\texttt{U+0041}, Я --- \texttt{U+042F}. UTF-8 и UTF-16 задают, как
представлять их кодовыми единицами и затем байтами.
Допустимые скалярные значения лежат от \texttt{U+0000} до
\texttt{U+10FFFF}, кроме суррогатного диапазона
\texttt{U+D800}--\texttt{U+DFFF}.

\begin{itemize}
    \item \textbf{UTF-8:} от 1 до 4 байт на скалярное значение.
    Коды ASCII представлены одним байтом с тем же значением.
    \item \textbf{UTF-16:} одна или две 16-битные кодовые единицы,
    то есть 2 или 4 байта. Символы за пределами основной многоязычной
    плоскости требуют суррогатной пары. При записи в байты важен порядок
    байт: UTF-16LE или UTF-16BE.
\end{itemize}
\textit{Пример:} строка «AЯ» в UTF-8 имеет байты
\texttt{41 D0 AF} (3 байта), а в UTF-16LE без BOM ---
\texttt{41 00 2F 04} (4 байта).
Кодовая точка \texttt{U+1F600} в UTF-8 занимает 4 байта
\texttt{F0 9F 98 80}, в UTF-16 --- две единицы \texttt{D83D DE00}.

Один видимый знак может состоять из нескольких кодовых точек:
например, буква и комбинируемый знак ударения. Поэтому число байтов,
кодовых единиц, кодовых точек и видимых знаков не обязано совпадать.
В C функция \texttt{strlen} считает байты до нулевого байта,
а не видимые символы UTF-8. Тип \texttt{char} сам по себе не задаёт кодировку.

\subsection*{Переход к числам}
Для чисел тоже необходимо договориться о формате: ширине слова,
наличии знака и правиле декодирования. Далее сначала рассматриваем
целые числа, затем приближённое представление вещественных чисел.
